\chapter{\rm\bfseries Background and Related Work}
\label{ch:background}

\section{Cloth Grasping in Robotics}

Robotic manipulation of cloth and other deformable fabric is difficult
precisely because the object being manipulated lacks a fixed geometry:
unlike a rigid part, a cloth's shape is itself part of the manipulation
state, and it must typically be inferred from vision rather than assumed
from a physical model. Early work in this space focused on perceiving that
shape well enough to act on it. Kita et al.~\cite{kita2004deformablemodel}
fit a deformable mass-spring model of a held garment to stereo-camera
silhouette observations, using the best-matching simulated deformation to
estimate the position of an unheld part of the garment for a second
manipulator to grasp. This was an early demonstration that a
physics-informed shape prior, rather than raw image features, is what
makes cloth state estimation tractable. Maitin-Shepard et
al.~\cite{maitinshepard2010clothgrasp} took a complementary, model-light
approach: a multiple-view geometric-cue detector locates cloth corners
directly from silhouettes, in a manner robust to texture and colour variation, letting
a general-purpose PR2 manipulator pick and fold towels of varying material
and appearance without a specialised end-effector. Miller et
al.~\cite{miller2012geometric} generalised this line of work with a
polygonal-approximation geometric model that tracks a garment's silhouette
through an entire folding sequence, allowing one parameterised folding
strategy to transfer across previously unseen towels, pants, and shirts.
Li et al.~\cite{li2015regrasping} pushed the physical modelling further
still, using a predictive thin-shell simulation fit to RGB-D observations
to plan regrasp actions that unfold and flatten a crumpled garment
by selecting the grasp points expected to maximise resulting flatness.

More recent work has moved from hand-designed geometric priors toward
learned policies and richer sensing. Saxena and
Shibata~\cite{saxena2019garment} trained convolutional classifiers to
recognise garment type and locate hem/waistline grasp points directly from
images, as the perception front-end of a robotic dressing-assistance
pipeline. Wu et al.~\cite{wu2020learning} showed that a model-free visual
reinforcement-learning policy for cloth and rope manipulation can be
learned an order of magnitude faster than a jointly-learned pick-and-place
policy by conditioning only on the \emph{placing} action and selecting
picks that maximise the value of the resulting placement, entirely without
human demonstrations. Closest in spirit to the force-sensing emphasis of
this thesis, Tirumala et al.~\cite{tirumala2022singulate} mounted a tactile
sensor on a Franka gripper fingertip and trained a classifier on the
resulting tactile signal to determine how many layers of cloth had been
grasped, using that signal to re-grasp until exactly one or two layers were
singulated from a stack. This is direct evidence that certain cloth-grasp ambiguities
cannot be resolved by vision alone. For a comprehensive treatment of the field beyond these
representative examples, Longhini et
al.~\cite{longhini2025unfolding} provide a recent survey spanning
perception, dynamics modelling, benchmarking, and control specifically for
textile manipulation.

\section{Strategies for Acquiring the First Grasp}
\label{sec:first-grasp-strategies}

Because a flat cloth on a surface presents no protruding feature for a
conventional pinch grasp, the literature has developed two broad families
of workarounds: manipulating the \emph{environment or cloth state} until a
graspable feature exists, and replacing the \emph{end-effector} with a
mechanism that does not require one.

\subsection{Environment-Assisted and State-Changing Actions}

The first family keeps a standard gripper and instead spends actions
creating a graspable feature (a fold, wrinkle, or free edge) before
the pick. Ramisa et al.~\cite{ramisa2012wrinkled} detected and ranked
wrinkles in depth images of crumpled laundry, directing a standard
gripper to the most graspable bump on the cloth rather than to a flat
region. Qian et al.~\cite{qian2020clothregion} segmented cloth edges
and corners (as distinct from wrinkles and folds) in depth images and
coupled the segmentation to a sliding grasp policy that approaches along
the estimated grasp direction. At the dynamic extreme, Ha and
Song's FlingBot~\cite{ha2021flingbot} showed that a high-velocity fling
primitive can unfold a crumpled cloth in a few actions: environment
interaction used not to create a pinchable feature but to reset the cloth
to a flatter, more predictable state. These strategies are effective, but
each spends additional actions and sensing before the grasp itself, and
their reliability degrades exactly in the flat-cloth-on-surface
configuration this thesis targets, where there is no wrinkle to find and
no free edge to slide against.

\subsection{Specialised End-Effectors}

The second family redesigns the gripper so that a protruding feature is
unnecessary. Monkman's survey of grippers for fibrous
materials~\cite{monkman1995fibrous} catalogues the classical industrial
approaches, developed largely for garment manufacturing: intrusive
(needle/pin) grippers that pierce the top ply, and adhesive or
electrostatic surface grippers. Electroadhesion has since matured substantially;
Guo et al.~\cite{guo2020electroadhesion} review its use across robotics
and highlight its suitability for delicate, flexible materials at low
energy cost. Non-contact airflow methods lift fabric without touching it
at all: Ozcelik and Erzincanli's Bernoulli-principle
end-effector~\cite{ozcelik2002noncontact} suspends a garment ply on a
radial airflow film. More recently, variable-friction fingertips blur the
line between gripper and manipulation policy: Jiang et
al.~\cite{jiang2022variablefriction} switch fingertip surfaces between
high- and low-friction states to slide, pinch, and unfold fabric in-hand.
These end-effectors trade generality for capability: needles risk damaging
fine textiles, electroadhesion and airflow devices require dedicated
power/pneumatic infrastructure and are largely single-purpose, and all of
them abandon the ubiquitous two-jaw form factor that general-purpose
robot platforms (including the Franka Hand used here) provide.

\subsection{The Gap: the Gripper as an Experimental Variable}

Across both families, the gripper hardware is treated as \emph{fixed} and
the innovation budget is spent on perception, policy, or a task-specific
mechanism. Direct comparisons between compliant and rigid jaw-type
grippers exist; Li et al.~\cite{li2025scopingreview} systematically
review soft--rigid hybrid grippers against industrial parallel-jaw
grippers across payload, adaptability, and precision metrics. Such
comparisons, however, are conducted on rigid or quasi-rigid workpieces,
not on cloth, and report capability envelopes rather than parametric contact
behaviour.

The closest precedent to this thesis is the household-clothing benchmark
of Clark et al.~\cite{clark2023household}, who evaluated several
commercial end-effectors (including finray-style compliant fingers
mounted on a Franka Emika Hand, alongside the standard Franka Hand,
a Robotiq 2F-140, and the OpenHand Model T42) on four cloth-specific
benchmarks: edge-drag accuracy, edge-grasp resilience, crumpled-object
encapsulation, and fold generation. Their results establish empirically
that end-effector selection materially changes cloth-manipulation
outcomes, and their inclusion of finray fingers on the same robot
platform used here makes the comparison direct. The evaluation, however,
is scored at the \emph{task} level: each benchmark reports success or
accuracy of a completed manoeuvre, with the approach pose fixed per task.
What it does not measure is how the grasp itself degrades as the
approach conditions vary: the continuous dependence of contact force
and grasp outcome on press depth and gripper tilt that any deployed
system experiences as positioning error.

None of the work above, therefore, addresses the specific question this
thesis investigates: cloth-manipulation literature treats grasp
\emph{success} as an outcome to be predicted or achieved at a fixed
operating point, not as a function of a continuous physical parameter
space (press depth, gripper tilt) whose boundary can be mapped and
compared across gripper geometries.
The benchmark in this thesis is deliberately narrower in manipulation scope
than folding or dressing pipelines, but correspondingly more quantitative
about the mechanics of the grasp itself: it keeps the two-jaw form factor
fixed, varies only the finger design $(t, \alpha)$ (defined in Figure~\ref{fig:seat-mount-params}), and maps the resulting
contact-force landscape under the positional and angular uncertainty
$(d, \theta_x, \theta_y)$ that any deployed system must tolerate.

\section{Compliant Gripper Designs}

The finray-effect finger design used in this thesis originates in a patent
by Kniese~\cite{kniese2000loadcarrying}, who observed that a structural
element built from two flexible longitudinal members joined by hinged
cross-struts deflects \emph{toward} an applied lateral load rather than
away from it, the counter-intuitive ``bends toward the pressure''
behaviour later commercialised by Festo as the Fin Ray Effect\textregistered.
Festo's own current technical documentation for its
DHAS gripper-finger product line~\cite{festo2018dhas} describes this
structure explicitly as two flexible bands joined by cross-ribs at flex
hinges, and reports measured retention-force, lateral-force, and
indentation-depth curves as functions of gripping force and workpiece
diameter. These are the same three quantities (force, indentation depth,
contact geometry) that this thesis's depth-sweep experiments characterise directly
on 3D-printed finray fingers rather than the commercial product.

Academic characterisation of the finray design has followed two
directions: analytic modelling and empirical/manufacturing studies. Shan
and Birglen~\cite{shan2020modeling} developed an analytic mechanical model
of finray fingers to predict grasp force and finger deformation from
geometric parameters, addressing the absence of a convenient design tool
for this class of compliant finger. Their model, like the empirical
studies that follow it, addresses \emph{lateral} grasping contact; no
published model, to the author's knowledge, predicts the axial
tip-compression response (the rise--peak--valley--recovery profile)
characterised empirically in this thesis. On the empirical side, Crooks et
al.~\cite{crooks2016finray} modified the finray design by angling its
internal crossbeams to introduce a directional bending bias, fabricated by
multi-material 3D printing, and reported an approximately 40\% increase in
holding capacity over a traditional finray finger from this geometric
change alone. Elgeneidy et al.~\cite{elgeneidy2019characterising}
characterised fully 3D-printed finray fingers made from flexible
thermoplastic polyurethane (TPU) filament, varying crossbeam rib angle and
spacing, and found that at large tip displacements the ribs contact one
another (\emph{layer jamming}), substantially increasing force generation
and producing a tunable-stiffness effect. Most directly relevant to this
thesis's use of TPU-printed fingers, Lang et
al.~\cite{lang2025softrobotics} 3D-printed finray fingers in both TPU 60A
and TPU 95A and tensile-tested them: the softer 60A durometer gave the
better force/compliance balance and left delicate fruit undamaged, while
95A reached substantially higher forces but visibly bruised the fruit.
That result is direct prior evidence that finger material hardness, independent of
finger geometry, is a first-order design lever for a compliant gripper, a
finding this thesis's own manufacturing-variability results
(Section~\ref{sec:manufacturing-results}) are consistent with.

The operational-envelope argument for compliant hands (that
compliance buys tolerance to positioning error) was established
quantitatively by Dollar and Howe's SDM Hand~\cite{dollar2010sdm}: an
underactuated hand with compliant polymer fingers that grasped a 48\,mm
cylinder at positions offset by more than 100\% of the object size in
each direction, while keeping acquisition contact forces low, where a
rigid gripper would require millimetre-level placement. Their
successful-grasp-range maps over an $xy$ offset grid are the direct
methodological ancestor of the depth--tilt working-window maps produced
in this thesis, transposed from rigid household objects to cloth and
from lateral offset to press depth and tilt. Closest of all in spirit,
Teeple et al.~\cite{teeple2022multidimensional} decomposed the benefit of
a pneumatic soft gripper's compliance into per-axis contributions when
handling thin, flexible materials: vertical compliance was shown to buy
robustness to large vertical positioning uncertainty, measured
head-to-head against a rigid gripper, while lateral and rotational
compliance lengthen the response window and cap the tensile force during
snags. Their vertical-uncertainty comparison is the single-gripper-pair
analogue of this thesis's central question; the benchmark here extends it
across a family of finger designs, adds the tilt dimension, and grounds
it on cloth with a fully passive compliant finger.

Compliant gripping is not achieved solely through finray-style flexure.
Granular jamming is the other major compliant-gripping paradigm in the
literature: Brown et al.~\cite{brown2010universal} demonstrated that a
latex membrane filled with granular material conforms to an object's
shape under positive pressure and rigidifies under vacuum, enabling
passive universal grasping without dedicated fingers at all, and Amend et
al.~\cite{amend2012positivepressure} extended this with a combined
positive/negative pressure actuation cycle for faster, more reliable
grasp-and-release cycles. A third paradigm, pneumatic elastomer actuation,
underlies commercial soft grippers such as those from Soft Robotics Inc.;
the technology traces to pneumatic-network (PneuNet) elastomer actuators
developed by Ilievski et al.~\cite{ilievski2011softrobotics}, and the
commercial mGrip product line~\cite{softroboticsinc_mgrip} demonstrates
that compliant gripping has moved beyond academic prototypes into
food-handling industrial deployment. This thesis's finray variants and
rigid parallel-jaw comparator sit within this broader design space as one
specific, geometrically simple point in it, chosen because (unlike jamming
or pneumatic actuation) a finray finger requires no additional actuation
hardware beyond the gripper's own jaw motion.

\section{Force-Controlled Grasping}

The theoretical basis for treating contact force, rather than position
alone, as a control variable during manipulation is Hogan's impedance
control framework~\cite{hogan1985impedancePartI,hogan1985impedancePartII,hogan1985impedancePartIII}:
neither pure position control nor pure force control is adequate for a
task involving contact with an environment of unknown or varying
mechanical properties, so the manipulator should instead regulate the
dynamic relationship --- the impedance --- between its motion and the
interaction force at its end-point. This thesis's trial procedure, which
presses the gripper to a commanded depth and then reports the resulting
\emph{stationary} contact force $F_z$ rather than commanding force
directly, is an impedance-control-style measurement: depth is the
commanded quantity, force is the observed response, and the relationship
between them (Chapter~\ref{ch:results}) is exactly the impedance curve
Hogan's framework predicts should characterise contact with an unknown
environment (here, the cloth-covered rigid table).

On the Franka FR3 platform specifically, libfranka exposes two
real-time command modes: external torque-level control (including a
Cartesian impedance controller distributed as a reference
implementation) and Cartesian pose motion generation, in which the
client streams pose targets at the 1\,kHz servo rate and tracking is
delegated to the arm's internal impedance
controller~\cite{frankarobotics_libfranka}. The tilt-aware pose
targeting used in this thesis's trial procedure
(Section~\ref{sec:tilt-pose}) uses the latter mode: the in-house
\texttt{fr3\_pose\_controller} commands Cartesian poses through
libfranka's pose interface, so the contact interaction is mediated by
the robot's internal impedance law rather than a user-tuned stiffness.

Contact force in this thesis is measured through the Franka FCI's own
external-wrench estimate; a second, independent channel using an array
of Tekscan FlexiForce resistive force sensors was provisioned, but a
repeatable calibration could not be established in time and the
cross-check is carried as future work (Chapter~\ref{ch:conclusion}). FlexiForce is a thin-film piezoresistive sensor
marketed explicitly for grip-force measurement in robotic and surgical
applications~\cite{tekscan_flexiforce_gripforce}, and has prior academic
precedent for exactly this use: Sharan and
Onwubolu~\cite{sharan2015flexiforce} instrumented a two-finger
pick-and-place gripper with a FlexiForce sensor to measure grip force
directly against known payload weight. An alternative tactile-sensing
approach not used in this thesis but relevant for future work is
vision-based tactile sensing: Yuan et al.'s
GelSight~\cite{yuan2017gelsight} sensor measures the deformation of a soft
elastomer contact surface via an internal camera, recovering high-resolution
contact geometry from which force and slip can be inferred, in contrast to
FlexiForce's direct (but lower-resolution) force read-out.

Finally, the trial procedure's touch-off step (descending the gripper
under position control until a chosen contact criterion is met, before
switching to the pressed-depth grasp) follows a long lineage of
contact-detection strategies in manipulation. Whitney's classical analysis
of quasi-static assembly~\cite{whitney1982quasistatic} derives the
geometric and force-equilibrium conditions under which measured contact
force can be used to detect and guide compliant insertion, the analytical
basis for the Remote Center Compliance device and, more generally, for
using force as a contact-detection signal rather than only as an outcome
measure. De Luca et al.~\cite{deluca2006collision} showed that contact can
be detected from proprioceptive joint measurements alone, without a
dedicated force/torque sensor, via a generalised-momentum residual, a
capability this thesis's FCI-based touch-off implicitly relies on when it
uses the controller's own external-wrench estimate rather than an external
load cell.

\section{The Franka FR3 Platform}
\label{sec:franka-platform-lit}

The Franka Research 3 (FR3) is a 7-degree-of-freedom, torque-controlled
collaborative arm with seven integrated joint torque sensors, a 3\,kg
payload, and position repeatability under $\pm 0.1$\,mm, accessed in this
thesis through the Franka Control Interface
(FCI)~\cite{franka2026fr3spec,frankarobotics2026fcidocs}, a real-time
control layer providing torque, position, and velocity access in both
joint and Cartesian space at up to 1\,kHz. This thesis controls the arm
through \texttt{franka\_ros2}~\cite{frankarobotics2026franka_ros2}, the official
ROS\,2 integration built on \texttt{libfranka} and \texttt{ros2\_control}.

The FR3's predecessor, the Franka Emika Panda, is documented by its
manufacturer as a widely adopted reference platform for robotics research
and education~\cite{haddadin2022franka}, and a more recent survey by the
same author quantifies its continued adoption as a de facto standard
platform across manipulation, control, and human-robot-interaction
research~\cite{haddadin2024standard}. That widespread adoption has
produced independent characterisation studies directly relevant to this
thesis's instrumentation choices. Most pertinently, Petrea et
al.~\cite{petrea2021interaction} empirically evaluated the accuracy of the
Panda's built-in filtered external-force estimate against ground-truth
load-cell measurements during controlled contact. This is the same FCI
external-wrench signal this thesis uses as its Contact Force channel, so
their accuracy findings directly bound how much that channel can be
trusted, and motivate the independent force-sensor cross-check
identified as future work.
Gaz et al.~\cite{gaz2019dynamic} separately derived a full dynamic
(inertial and friction) parameter model of the Panda's arm via
penalty-based optimisation, a model since reused widely enough to be
treated as a de facto standard dynamic model of the platform. Beyond force
sensing specifically, the Franka Panda/FR3's broad adoption as a
manipulation-research testbed is illustrated by work such as Chi et
al.'s~\cite{chi2023diffusion} diffusion-policy visuomotor control
framework, evaluated on contact-rich real-world manipulation tasks using a
Franka Panda as one of its physical platforms. This adoption is directly
relevant to the control mode chosen here: learned visuomotor policies of
this kind emit streams of end-effector pose targets, the same command
interface through which this thesis's trial procedure drives the arm
(Section~\ref{sec:tilt-pose}). The benchmark is therefore conducted
through the interface a learned or simulation-trained policy would use in
deployment, with sim-to-real transfer as the motivating trajectory: a
compliant finger that tolerates centimetre-scale press-depth error and
degree-scale tilt error widens the placement-error budget available to a
policy trained in simulation, and the depth--tilt maps produced in this
thesis quantify that budget as a function of gripper design.

\section{Position and Summary}

This thesis sits at the intersection of three threads: cloth-handling
system design, finray-effect gripper characterisation, and
force-controlled manipulation on collaborative robots. Prior work has
independently established each thread: cloth manipulation research has
largely optimised for task-level success (folding, singulating, dressing)
under vision-driven or learned control; compliant-gripper research has
characterised finray and jamming designs primarily through force/holding-capacity
metrics on rigid or quasi-rigid test objects; and impedance-control theory
and its Franka FCI implementation provide the tools to regulate and
measure contact force but have not, in the reviewed literature, been used
to map how that force varies over a continuous depth--tilt parameter space
for a fragile, non-rigid workpiece. The contribution of this thesis is a
unified experimental treatment of all three on a single platform with
consistent instrumentation, enabling direct quantitative comparison
between rigid and compliant gripper response under deliberately perturbed
conditions.
